\begin{table}[t]
\centering
\caption{Observational Image-Space Diagnostics on SLACS HST Lenses}
\label{tab:real_data}
\small
\begin{tabular}{l|cc|cccc|c}
\toprule
\textbf{Lens} & $z_l$ & $\sigma_v$ & \textbf{NRMSE}$\downarrow$ & \textbf{SSIM}$\uparrow$ & \textbf{Corr.}$\uparrow$ & \textbf{Flux Ratio} & \textbf{Pass} \\
\midrule
  J0946+1006 & 0.222 & 263 & 0.08953 & 0.9868 & 0.931 & 0.952 & $\checkmark$ \\
  J1250+0523 & 0.232 & 252 & 0.10594 & 0.9809 & 0.928 & 0.967 & $\checkmark$ \\
  J1402+6321 & 0.205 & 267 & 0.10310 & 0.9882 & 0.899 & 0.965 & $\checkmark$ \\
  J0252+0039 & 0.280 & 164 & 0.11232 & 0.9881 & 0.889 & 0.939 & $\checkmark$ \\
  J0037-0942 & 0.195 & 279 & 0.09252 & 0.9831 & 0.935 & 0.946 & $\checkmark$ \\
\bottomrule
\end{tabular}
\vspace{1mm}
\parbox{\columnwidth}{\footnotesize Velocity dispersions $\sigma_v$ (km/s) and lens redshifts from Bolton et al. (2008). The lens mass model is fixed from the published SLACS parameters, then a foreground-suppressed annular HST image is fit with a PSF-convolved lensed S\'ersic source model. The flux-ratio column reports annular model/observation flux consistency.}
\end{table}
